\documentclass[11pt, a4paper]{article}

\usepackage[T1]{fontenc}
\usepackage[utf8]{inputenc}
\usepackage[english]{babel}

% times new roman and matching math font
\usepackage{libertine}
\usepackage{newtxmath}

\usepackage{booktabs}
\usepackage[svgnames]{xcolor}
\usepackage{graphicx}
\usepackage{hyperref}
\usepackage{parskip}
\usepackage{url}

\setlength{\oddsidemargin}{0cm}
\setlength{\evensidemargin}{0cm}
\setlength{\topmargin}{-1cm}
\setlength{\textheight}{23cm}
\setlength{\textwidth}{16cm}

%\setlength{\parindent}{0pt}

\pagestyle{headings}

\newcommand{\bat}{{\sc BAT}}
\newcommand{\BAT}{\bat}
\newcommand{\Root}{{\sc Root}}

\newcommand{\versionno}{1.1.0-DEV}

\newcommand{\version}{version~\versionno}
\newcommand{\Version}{Version~\versionno}

\newcommand{\code}[1]{\texttt{#1}}

% comments
\newcommand{\dan}[1]{{\color{red}{Dan: #1}}}
\newcommand{\fred}[1]{{\color{blue}{Fred: #1}}}

% references to ...
\def \refeq#1{(\ref{eq:#1})}
\def \refsec#1{Section~\ref{sec:#1}}
\def \reffig#1{Figure~\ref{fig:#1}}
\def \reftab#1{Table~\ref{tab:#1}}

% syntax highlighting
\usepackage{listings}
\definecolor{mygray}{rgb}{0.4,0.4,0.4}
\definecolor{mygreen}{rgb}{0,0.8,0.6}
\definecolor{myorange}{rgb}{1.0,0.4,0}
\lstset { %
  backgroundcolor=\color{black!5}, % set backgroundcolor
  basicstyle=\ttfamily,% basic font setting
  % keywordstyle=\color{Magenta},
}
\lstnewenvironment{cxxcode}
{\lstset{
    language=c++,
    gobble=0,
    commentstyle=\itshape\color{purple!80},
    breaklines=true
    basicstyle=\sffamily\color{black},
    frame=single,
    % identifierstyle=\color{blue},
    keywordstyle=\color{mygreen},
    numberstyle=\tiny\color{mygray},
    showspaces=false,
    showstringspaces=false,
    stringstyle=\color{myorange},
    tabsize=4
  }
}
{}

%--------------------------------------------------------

\begin{document}

% --------------------------------------------------------
% title
% --------------------------------------------------------

\thispagestyle{empty}

\begin{figure}
\includegraphics[scale=0.25]{bat}
\end{figure}

\vspace*{1cm}

\begin{center}


{\Large An introduction to \bat}
\\

\vspace{1cm}

{\large \bat\ \version}

\end{center}

\thispagestyle{empty}

\vfill

\begin{center}
\today
\end{center}

\pagebreak

% --------------------------------------------------------
% table of contents
% --------------------------------------------------------

\thispagestyle{empty}

\enlargethispage{2cm}

\tableofcontents

\section{Overview}

\fred{what is this for, web page, reference guide, github, teaser}

The purpose of this document is to introduce the Bayesian Analysis
Toolkit (BAT). For a quick tutorial covering the basic features, jump
directly to \refsec{basics}.

\subsection*{Structure of this document}

We give a brief introduction of the Bayesian approach to statistics in
\refsec{Bayes} and describe the core algorithm of \bat{}, Markov chain
Monte Carlo, in \refsec{MCMC}. Further sections describe other parts
of \bat{} and advanced features.

\subsection*{Further help}

Please visit our home page at \url{http://mpp.mpg.de/bat} for an
overview of documentation and activities around \bat{}.

This document is meant to is quickly enable to you to find your way
around \bat{}. We describe the common workflow and explain how to
solve some not-so-common tasks. To keep the size of this document
manageable, we refer you to the
\href{http://mpp.mpg.de/bat/docs/refman/latest/}{online reference
  guide} in which all classes, methods etc. are documented.

If you find a problem with the code, the preferred method is to create
a new issue at \url{https://github.com/bat/bat/issues}. For general
questions or comments, contact the developers at
\href{mailto:bat@mpp.mpg.de}{bat@mpp.mpg.de}.

\section{Basic BAT} \label{sec:basics}

Goal of this section is that if you can read and test as you go along, and end up with functioning code.

\subsection{Installation instructions}

Detailed instructions for downloading and installing BAT are in the
\verb|INSTALL.md| file that comes with the \bat\ source distribution. The
latest version is available online at
\url{https://github.com/bat/bat/blob/master/INSTALL.md}.

\subsection{Non-code parts of BAT}
\dan

bat-project, bat-config
\begin{cxxcode}
int main() {
  std::string s = "Hello";
  return 0;
}
\end{cxxcode}

\subsection{Absolutely necessary requirements of your model}
\dan

Parameters, LogLikelihood, some form of setting a prior.

\subsection{Running the analysis}
\dan

Functions to call. How to output. Basic tuning. How to interpret runtime log output \& trouble shooting (nonconvergence, poor efficiencies). How to interpret summary log output.

\subsection{Plots}
\dan

Interpreting the basic output: marginals; parameters; knowledge update; correlation matrix.

Options for the different output types.

\section{Intro to Bayesian Statistics (for our purposes)} \label{sec:Bayes}
\fred

\section{Intro to Markov Chain Monte Carlo (for our purposes)} \label{sec:MCMC}
\fred

R value.
\section{Non-MCMC parts of BAT}

\subsubsection{Integration}
\fred

\subsubsection{Optimization}
\fred

\section{Accompanying Models}

\subsection{Using BAT's native data structures}

\subsection{base}

\subsection{mtf}


\section{Advanced BAT}

\subsection{ModelManager for Model Comparison}
\fred{can work on same data set but not required}


\subsection{Defining your own factorized prior}
\dan

\subsection{Sharing / Loading samples}
\dan{describe that files can be checked during prerun if \#136 is merged}

\subsection{Performance / Threading / User's attention to thread-safety}

tips.

\section{Examples}

\section{Validation}

Proof we've checked BAT works.

\section{Troubleshooting or FAQ}

\section{Further reading}


\end{document}
